On introduit dans un bécher un volume $V_{1}=\qty{500}{\cubic\cm}$ d'eau
distillée et on lui ajoute un volume $V_{2}=\qty{1}{\cubic\cm}$ d'une solution
d'acide éthanoïque pure.
On mesure de le pH du mélange à l'aide d'un pH mètre et on obtient~:
$\mathrm{pH}=3,1$.

\textbf{Données~:}
\begin{itemize}
    \item La densité de l'acide éthanoïque~: $d=1,05$
    \item La masse volumique de l'eau~:
          $\rho_{\text{eau}}=\qty{1}{\gram\per\cubic\cm}$
    \item La masse molaire de la molécule d'acide éthanoïque~:
          $M(\ce{CH3COOH})=\qty{60}{\gram\per\mole}$.
\end{itemize}

\begin{enumerate}
    \item Écrire la réaction chimique de l'acide éthanoïque avec l'eau.
    \item Déterminer la quantité de matière initiale de l'acide éthanoïque.
    \item Dresser le tableau d'avancement de la réaction puis déterminer la
          valeur de l'avancement maximal.
    \item Déterminer la valeur de l'avancement final. Quelle est votre
          conclusion.
    \item Calculer le taux d'avancement final de la réaction.
\end{enumerate}

